\chapter[The tensor mass gap]{Vanishing of the Tensor Mass Gap}
\label{ch:massless-tensor}

\paragraph{Learning outcomes.}
After this chapter the reader should be able to:
\begin{enumerate}
  \item write the Fierz--Pauli quadratic action for \(T_{\mu\nu}\)
        and the linearized Einstein operator that remains when
        \(m_T=0\);
  \item substitute a compactly supported tensor variation into the
        entanglement Ward identity and see where \(m_T^2\) appears;
  \item state Theorem~\ref{thm:mT} with its hypotheses, and say which
        hypotheses are still assumptions of the framework;
  \item count the two remaining TT polarizations;
  \item list what must still be done for the conformal and vector
        blocks before the infrared theory is pure GR.
\end{enumerate}

The candidate graviton now has a name, a quadratic action, and a
Ward identity. This chapter puts the last of those into the second
and reads off \(m_T=0\).

\section{Strategy}
\label{sec:ch7-strategy}

On a Minkowski condensate the tensor block is a linear field
\(T_{\mu\nu}\) with the algebraic symmetries~\eqref{eq:STF}. Its
quadratic action has the generic structure~\eqref{eq:ST-generic}.
The entanglement Ward identity, Proposition~\ref{prop:ward}, must
hold for every allowed variation supported on that block.

The geometric side of the identity, \(\delta A\), depends on
\(T_{\mu\nu}\) only through the linearized metric
\begin{equation}
  \label{eq:h-is-T}
  h_{\mu\nu}=T_{\mu\nu}+\tfrac14\eta_{\mu\nu}\psi,
\end{equation}
and, in a gauge where the conformal piece is held fixed by the
fixed-volume constraint of Assumption~\ref{ass:EE}, through
\(T_{\mu\nu}\) alone. The modular side,
\(\delta\langle T_{\mu\nu}^{\mathrm{op}}\rangle\), is the metric
stress of the quadratic tensor action. If that action contains a
mass term, the modular side contains a piece proportional to
\(m_T^2 T_{\mu\nu}\). The geometric side does not. Compatibility for
arbitrary compactly supported \(T_{\mu\nu}\) forces \(m_T^2=0\).

That is the argument. The rest of the chapter writes it so that each
step can be rejected independently.

\section{Fierz--Pauli action for the tensor block}
\label{sec:ch7-FP}

The unique two-derivative quadratic action for a massive symmetric
tensor that is ghost-free about Minkowski spacetime is the
Fierz--Pauli action. After the algebraic trace has been projected
out, so that \(T^\mu{}_\mu=0\) identically, it reduces to
\begin{equation}
  \label{eq:FP}
  \begin{aligned}
    S_{\mathrm{FP}}[T]
    =\int\mathrm{d}^4x
     \Bigl(
       &\tfrac12\partial_\rho T_{\mu\nu}\partial^\rho T^{\mu\nu}
        -\partial^\rho T_{\rho\mu}\partial_\sigma T^{\sigma\mu}\\
       &-\tfrac12 m_T^2 T_{\mu\nu}T^{\mu\nu}
     \Bigr).
  \end{aligned}
\end{equation}
The derivative structure is required. A generic mass term
\(m^2(T_{\mu\nu}T^{\mu\nu}-\alpha T^2)\) with \(\alpha\neq 1\)
(before the trace projection) propagates a ghost. With the trace
already zero the remaining danger is a wrong-sign kinetic term,
which~\eqref{eq:FP} does not have.

The Euler--Lagrange equation of~\eqref{eq:FP} is
\begin{equation}
  \label{eq:FP-eom}
  \mathcal{E}_{\mu\nu}[T]-m_T^2 T_{\mu\nu}=0,
\end{equation}
where
\begin{equation}
  \label{eq:einstein-op}
  \begin{aligned}
    \mathcal{E}_{\mu\nu}[T]
    &\equiv
    \square T_{\mu\nu}
    -\partial_\mu\partial^\rho T_{\rho\nu}
    -\partial_\nu\partial^\rho T_{\rho\mu}\\
    &\qquad
    +\eta_{\mu\nu}\partial^\rho\partial^\sigma T_{\rho\sigma}
  \end{aligned}
\end{equation}
is the linearized Einstein operator on a traceless tensor about
Minkowski spacetime. (The linearized Einstein tensor on a general
\(h_{\mu\nu}\) contains additional trace pieces; they are absent
here because \(T^\mu{}_\mu=0\).)

When \(m_T=0\), \eqref{eq:FP} is invariant under
\begin{equation}
  \label{eq:FP-gauge}
  \delta_\xi T_{\mu\nu}
  =\partial_\mu\xi_\nu+\partial_\nu\xi_\mu
   -\tfrac12\eta_{\mu\nu}\partial\cdot\xi
\end{equation}
with \(\partial\cdot\xi\) arranged so that the trace remains zero.
When \(m_T\neq 0\), that invariance is absent and the on-shell count
is five degrees of freedom.

\section{Modular stress of the tensor block}
\label{sec:ch7-stress}

The operator \(T_{\mu\nu}^{\mathrm{op}}\) in
Proposition~\ref{prop:ward} is defined by the same rule
as~\eqref{eq:T-def}, applied to \(S_{\mathrm{FP}}\) treated as a
matter action on the background metric that defines the diamond.
To first order about Minkowski spacetime,
\begin{equation}
  \label{eq:Top-FP}
  \delta\langle T_{\mu\nu}^{\mathrm{op}}\rangle
  =\mathcal{E}_{\mu\nu}[T]-m_T^2 T_{\mu\nu}
   +\partial^\rho\Sigma_{\rho\mu\nu},
\end{equation}
where \(\Sigma_{\rho\mu\nu}\) is a total-derivative superpotential
coming from the fact that the canonical and metric stress tensors
differ by an improvement. On a diamond the superpotential
contributes only at the edge \(\mathcal{B}\). That edge contribution
is absorbed into the area term of the first law, which is already
present in~\eqref{eq:ward}. What remains in the bulk integral is
the improved stress
\begin{equation}
  \label{eq:Top-improved}
  \delta\langle T_{\mu\nu}^{\mathrm{op}}\rangle_{\mathrm{bulk}}
  =\mathcal{E}_{\mu\nu}[T]-m_T^2 T_{\mu\nu}.
\end{equation}

\section{Area variation induced by \(T_{\mu\nu}\)}
\label{sec:ch7-area}

The edge \(\mathcal{B}\) of a small diamond of radius \(\ell\) is a
2-sphere of area \(4\pi\ell^2\) in the Minkowski background. A
metric perturbation \(h_{\mu\nu}=T_{\mu\nu}\) changes that area by
\begin{equation}
  \label{eq:deltaA}
  \delta A
  =\int_{\mathcal{B}}
   \mathrm{d}A\,
   \tfrac12\sigma^{\mu\nu}T_{\mu\nu}
   +O(\ell^3\|T\|,\ell\|T\|^2),
\end{equation}
where \(\sigma^{\mu\nu}\) is the induced metric on \(\mathcal{B}\).
Jacobson's geometric identity converts the same variation, after
integration against the conformal Killing vector and use of the
fixed-volume constraint, into
\begin{equation}
  \label{eq:deltaA-G}
  \frac{\kappa}{8\pi\Geff}\,\delta A
  =
  -\int_{\Sigma\cap D}
   \mathrm{d}\Sigma^\mu\,\xi^\nu
   \,\mathcal{E}_{\mu\nu}[T]
   +O(\ell^3).
\end{equation}
The right-hand side contains the Einstein operator and does
\emph{not} contain \(m_T\). The area of an emergent 2-surface knows
about the metric perturbation. It does not know about a potential
term in the fluctuation action.

Equation~\eqref{eq:deltaA-G} is imported from the geometry of small
diamonds, not re-derived here. Its hypotheses are: Riemann normal
coordinates to the order kept in~\eqref{eq:RNC}, fixed volume, and
the identification of the area prefactor with \(\Geff\). The latter
remains conditional on the specified sector calculation in Problem~F;
Chapter~8 gives only the truncated heat-kernel identity.

\section{Substitution into the Ward identity}
\label{sec:ch7-substitution}

Insert~\eqref{eq:Top-improved} and~\eqref{eq:deltaA-G} into
\eqref{eq:ward}. The Einstein-operator terms cancel, and what
remains is
\begin{equation}
  \label{eq:mass-identity}
  \int_{\Sigma\cap D}
  \mathrm{d}\Sigma^\mu\,\xi^\nu
  \bigl(-m_T^2 T_{\mu\nu}\bigr)
  +O(\ell^3)
  =0.
\end{equation}
The conformal Killing vector \(\xi^\nu\) of a small diamond is
timelike in the interior and vanishes on the edge. Its leading
expression in Riemann normal coordinates is
\begin{equation}
  \label{eq:xi-RNC}
  \xi^\mu
  =\frac{\kappa}{2\ell}
   \bigl((\ell^2-t^2-r^2)\partial_t^\mu+2t\,x^\mu\bigr)
   +O(R_{\alpha\beta\gamma\delta}\ell^3).
\end{equation}

\begin{lemma}[Arbitrary tensor sources]
  \label{lem:arbitrary}
  Compactly supported, smooth, traceless, symmetric variations
  \(T_{\mu\nu}\) are allowed in the sense of
  Definition~\ref{def:allowed}. In particular they may be chosen
  independently in each small diamond, up to the fixed-volume
  constraint already used.
\end{lemma}

\begin{proof}
  The tensor block was split off algebraically in Chapter~5. Gauge
  fixing of diffeomorphisms and local Lorentz transformations
  removes \(\psi\)-mixing and \(\omega\)-mixing at this order but
  does not constrain the nine components of \(T_{\mu\nu}\) beyond
  smoothness and compact support. Fixed volume constrains the
  trace sector, which has already been projected out.
\end{proof}

\begin{lemma}[Integral identity implies point identity]
  \label{lem:point}
  If~\eqref{eq:mass-identity} holds for every small diamond and
  every allowed \(T_{\mu\nu}\), then
  \begin{equation}
    \label{eq:point}
    m_T^2 T_{\mu\nu}(x)=0
  \end{equation}
  at every event \(x\) in the condensed phase.
\end{lemma}

\begin{proof}
  Suppose \(m_T^2 T_{\mu\nu}\) is not identically zero. Then there
  is an event \(x_0\) and a pair of indices at which the tensor
  is nonzero. Choose a diamond centered at \(x_0\) with \(\ell\)
  small enough that \(T_{\mu\nu}\) does not change sign in the
  relevant component on \(\Sigma\cap D\), and choose the time
  orientation so that \(\xi^\nu\) has a nonvanishing contraction
  with that component. The integrand of~\eqref{eq:mass-identity} is
  then of definite sign on a set of positive measure, and the
  \(O(\ell^3)\) remainder cannot cancel it as \(\ell\to 0\).
\end{proof}

\section{The vanishing theorem}
\label{sec:ch7-theorem}

\begin{theorem}[Vanishing of the tensor mass gap]
  \label{thm:mT}
  Assume:
  \begin{enumerate}
    \item the microscopic setup of Chapter~3, including
          Assumptions~\ref{ass:aux-metric}--\ref{ass:dictionary};
    \item the condensed frame of Chapter~4, including
          Assumptions~\ref{ass:nondeg} and~\ref{ass:lorentzian};
    \item the algebraic split of Chapter~5 and
          Assumption~\ref{ass:locked-block};
    \item entanglement equilibrium, Assumption~\ref{ass:EE}, in the
          form of Proposition~\ref{prop:ward};
    \item the Fierz--Pauli quadratic action~\eqref{eq:FP} as the
          two-derivative tensor block;
    \item the geometric identity~\eqref{eq:deltaA-G} for small
          diamonds.
  \end{enumerate}
  Then \(m_T^2=0\) throughout the condensed phase. The quadratic
  tensor action reduces to the linearized Einstein--Hilbert action
  on traceless perturbations,
  \begin{equation}
    \label{eq:EH-lin}
    S_{\mathrm{EH}}^{(2)}[T]
    =\int\mathrm{d}^4x\,
     \tfrac12 T^{\mu\nu}\mathcal{E}_{\mu\nu}[T],
  \end{equation}
  and is invariant under~\eqref{eq:FP-gauge}.
\end{theorem}

\begin{proof}
  Lemma~\ref{lem:point} gives \(m_T^2 T_{\mu\nu}=0\) at every event.
  The tensor block is not the zero field: linearized gravitational
  waves, if they exist, are particular solutions for \(T_{\mu\nu}\).
  Hence \(m_T^2=0\). Substituting into~\eqref{eq:FP} produces
  \eqref{eq:EH-lin} after an integration by parts, boundary terms
  being discarded for compactly supported variations. Gauge
  invariance~\eqref{eq:FP-gauge} is the standard invariance of the
  linearized Einstein operator on a traceless tensor.
\end{proof}

\begin{remark}[What the theorem does not hide]
  Hypothesis~4 is the physical input of NAMCG. Hypothesis~5 excludes
  non-Fierz--Pauli two-derivative structures and higher-derivative
  quadratic terms; the latter return in Chapter~9. Hypothesis~6
  borrows a geometric identity used in Jacobson's argument. The
  theorem is therefore a compatibility statement: a Fierz--Pauli
  tensor block cannot sit in an entanglement-stationary condensate
  unless it is massless. It is not a derivation of
  Assumption~\ref{ass:EE} from the path integral.
\end{remark}

\section{Two transverse-traceless polarizations}
\label{sec:ch7-TT}

With \(m_T=0\), the equation \(\mathcal{E}_{\mu\nu}[T]=0\) plus the
gauge invariance~\eqref{eq:FP-gauge} is linearized vacuum GR in the
traceless sector. Choose the de Donder gauge
\begin{equation}
  \label{eq:dedonder}
  \partial^\mu T_{\mu\nu}=0.
\end{equation}
Residual gauge transformations with \(\square\xi^\mu=0\) remain.
They remove four more components. Together with the four conditions
\eqref{eq:dedonder} and the four conditions implicit in taking the
trace to vanish and staying in the symmetric sector after gauge
fixing, one is left with two independent components. In a plane-
wave basis they are the helicities \(\pm 2\),
\begin{equation}
  \label{eq:TT-wave}
  T_{\mu\nu}(x)
  =\mathrm{Re}\bigl(
     \epsilon_{\mu\nu}^{\pm 2}\,
     e^{ik\cdot x}
   \bigr),
  \qquad
  k^2=0,
  \quad
  k^\mu\epsilon_{\mu\nu}=0,
  \quad
  \epsilon^\mu{}_\mu=0.
\end{equation}

This is the infrared graviton advertised in Chapter~1. It has been
obtained as a fluctuation of the condensate, not as a fundamental
metric field, contingent on Theorem~\ref{thm:mT}.

\section{The van Dam--Veltman--Zakharov warning}
\label{sec:ch7-vDVZ}

A massive spin-2 field with \(m_T\) small but nonzero does not
approach the massless theory smoothly in four dimensions. The
discontinuity is visible in the linearized deflection of light and
in the tensor-to-scalar ratio of a static source. Therefore it is
not enough to argue that \(m_T\) is ``small compared with
\(\LUV\).'' Entanglement equilibrium is being asked to set
\(m_T\) identically to zero in the infrared, not to suppress it by
a loop factor.

If a later calculation produced \(m_T\sim v\,e^{-\LUV/v}\) or any
other parametrically small but nonzero value, the theory would fail
Solar-System tests unless an additional Vainshtein mechanism were
demonstrated. No such mechanism is claimed here. Theorem~\ref{thm:mT}
is an identity-zero statement.

\section{Conformal and vector blocks}
\label{sec:ch7-other}

Theorem~\ref{thm:mT} is silent about \(\psi\) and \(A_\mu\).

\paragraph{Vector block.}
The internal gauge connection \(A_\mu\) is distinct from the
independent Lorentz connection \(\omega_\mu\) and from their diagonal
stabilizer combination. Its mass, if any, depends on kinetic
normalizations and the residual internal symmetry, neither of which
is specified here. If the internal gauge group has an unbroken
subgroup, the corresponding field is massless. Neither case is
claimed to produce a second graviton. Chapter~10 identifies the
broken-direction combination and the unitary-gauge removal of the
Goldstone orbit variables; it does not settle the diagonal
connection's dynamics (Problem~D).

\paragraph{Conformal block.}
The scalar \(\psi\) is more dangerous. If it remains massless it
is a Brans--Dicke field, and the infrared theory is scalar--tensor
gravity. Entanglement equilibrium at fixed volume was designed to
remove a conformal breathing from the comparison diamond; it does
not, by itself, gap \(\psi\) as a propagating excitation.

\begin{assumption}[Gapped conformal mode]
  \label{ass:gap-psi}
  The radial Higgs \(\sigma\) of the polar decomposition, and the
  conformal metric mode \(\psi\) mixed with it, acquire a mass
  \(m_\psi\sim\sqrt{\lambda}\,v\) from the potential~\eqref{eq:V}
  and decouple from the infrared tensor sector.
\end{assumption}

Assumption~\ref{ass:gap-psi} is the scalar-tensor ceiling of
Chapter~1, reappeared as a named hypothesis. If it fails, NAMCG
predicts a light scalar and is constrained by Cassini-type bounds
on \(\omega_{\mathrm{BD}}\). Chapter~17 records this as a
falsifier, not as a detail.

\section{Exercises}
\label{sec:ch7-exercises}

\begin{enumerate}
  \item \textbf{Fierz--Pauli ghost.}
    Restore a trace piece \(T^\mu{}_\mu=\psi\) and replace the mass
    term by \(m^2(T_{\mu\nu}T^{\mu\nu}-\alpha\psi^2)\). Show that
    \(\alpha\neq 1\) produces a sixth, ghostlike excitation.
    This is why~\eqref{eq:FP} is not an arbitrary massive action.
  \item \textbf{Gauge invariance at \(m_T=0\).}
    Verify that~\eqref{eq:FP-gauge} leaves~\eqref{eq:EH-lin}
    invariant. Show the mass term in~\eqref{eq:FP} is not
    invariant.
  \item \textbf{Einstein operator on a TT wave.}
    Insert~\eqref{eq:TT-wave} into~\eqref{eq:einstein-op} and
    recover \(k^2=0\).
  \item \textbf{Sign in the Ward identity.}
    Track the sign from~\eqref{eq:Top-improved} through
    \eqref{eq:deltaA-G} to~\eqref{eq:mass-identity}. A sign error
    here would force \(m_T^2=0\) by the same logic but would
    mismatch the Einstein equation in Chapter~11. Fix the sign
    against Jacobson's first law for ordinary matter.
  \item \textbf{Why compact support.}
    Suppose \(T_{\mu\nu}\) is a plane wave, not compactly supported.
    Where does the proof of Lemma~\ref{lem:point} fail, and how
    does one repair it by using wave packets?
  \item \textbf{vDVZ.}
    Look up the linearized deflection angle for Fierz--Pauli
    gravity with mass \(m_T\) and take \(m_T\to 0\). Compare with
    the GR deflection. What numerical factor appears?
  \item \textbf{Conformal leftover.}
    Linearize a Brans--Dicke action about Minkowski spacetime and
    identify the kinetic mixing of \(\psi\) with \(T_{\mu\nu}\).
    What bound on \(m_\psi\) follows from requiring the mixing to
    be negligible inside the Solar System?
\end{enumerate}

\section{Bridge to Chapter~8}
\label{sec:ch8-bridge-from-7}

The infrared tensor sector is now linearized general relativity,
contingent on Theorem~\ref{thm:mT} and Assumption~\ref{ass:gap-psi}.
The prefactor of that action is still a symbol, \(\Geff\). Chapter~8
records a conditional truncated heat-kernel identity; the physical
coefficient requires the specified sector calculation of Problem~F.
The coefficient is not evaluated here:
\begin{equation}
  \label{eq:Geff-next}
  C_R^L[\{\mathcal O_s,w_s,{\LUV}_s,\mu_s\},\mathcal S]
  \quad\text{is not yet matched to }G_{\mathrm{eff}}.
\end{equation}
The same expansion produces a quartic term \(c_0\LUV^4\), which
Chapter~8 will isolate and Chapter~12 will be required to tame.